\section{Introduction}
\label{sec:intro}

Most succinct arguments fix a single finite field and arithmetize every
operation into it. For bit-oriented computation --- hash functions, block
ciphers, the bitwise glue inside signature schemes --- this is a poor fit. A
$32$-bit $\XOR$, a rotation, an addition modulo $2^{32}$: each is a handful of
gates in hardware, but forced into a large prime field $\F_q$ it expands into
many constraints and many witness elements, because $\F_q$ has no native
notion of ``the $i$-th bit'' or ``the carry into position $j$''. The witness
inflates by one to two orders of magnitude, and since prover cost tracks
witness size, the encoding overhead dominates.

SHA-256 is the canonical stress test. Its compression function is built almost
entirely from three primitives on $32$-bit words --- bitwise $\XOR$, bitwise
$\AND$ (inside $\Ch$ and $\Maj$), and addition modulo $2^{32}$ --- together
with fixed rotations and shifts. A prover that treats a word as an opaque field
element must re-derive its bit structure, with range checks and bit
decompositions, at every step.

\paragraph{The all-$\Ftwo[X]$ approach.} The arithmetic SHA-256 actually
performs is over the integers modulo $2^{32}$, and the general Zinc\texttt{+}
pipeline \cite{zincplus} proves such computations by lifting them to $\Z[X]$.
But SHA-256 is \emph{overwhelmingly bitwise}, and for a bitwise computation the
productive proof-time algebra is characteristic two. We represent each
$\wbit$-bit word ($\wbit = 32$) as a \emph{bit-polynomial} in the cell ring
\[
  \cellring \;=\; \Ftwo[X]/(X^{\wbit}),
\]
the degree-$<\wbit$ polynomial over $\Ftwo$ whose coefficient of $X^b$ is the
word's $b$-th bit. In this representation the bitwise primitives are ring
operations rather than emulations:
\begin{itemize}[leftmargin=2em]
  \item \textbf{$\XOR$ is addition} in $\cellring$ ($1+1 = 0$ coefficientwise).
    An $\XOR$ of committed columns is a committed linear combination, never
    materialised --- free at the commitment layer.
  \item \textbf{Shifts and rotations are monomial multiplication} by $X^c$
    (in the shift quotient $\Fshift$ or rotation quotient $\Frot$), hence
    $\Ftwo$-linear; a shifted column is \emph{virtual}, reconstructed from its
    source rather than committed.
  \item \textbf{Modular addition is the Binius full-adder identity} in
    $\Fshift$, where multiplication by $X$ is a left shift and $X^{\wbit}=0$
    discards the overflow. The cell ring makes it cheap: the carry word is a
    free shift of committed data, so the prover commits only the \emph{one}
    carry-out bit per addition that the shift cannot supply.
  \item \textbf{$\AND$ is a Hadamard product} (hence so are $\Ch$ and $\Maj$):
    the coefficient-wise product of bit-polynomials, the single genuinely
    nonlinear primitive and the one the linear commitment cannot absorb.
\end{itemize}
The one price of working in characteristic two rather than over $\Z$ is modular
addition, which is not a ring operation; the third bullet is how it is paid.
The result is an AIR of SHA-256 in which \emph{every constraint is
$\Ftwo$-linear except a small fixed number of column Hadamard products}.

\paragraph{Three layers.} Turning this representation into a SNARK has three
layers, and the paper treats all three.
\emph{(1) Arithmetization} (\Cref{sec:arith}): the chained-compression AIR ---
the column model, the linear identities for the message schedule, the
$\Sigma/\sigma$ mixing, and modular addition, and the declaration of $\AND$ as
a column Hadamard relation.
\emph{(2) The binary-field PIOP} (\Cref{sec:piop,sec:hadamard}): the polynomial
interactive oracle reduction that proves the AIR. Because the trace already
lives in $\Ftwo[X]\subset\K[X]$, the prime-projection step of the general
Zinc\texttt{+} compiler is vacuous; what remains is an ideal check, an
\emph{evaluation projection} $\psia$ that substitutes $X = \alpha$ for a random
$\alpha\in\K = \mathrm{GF}(2^{128})$, and a sumcheck over $\K$. The nonlinear
Hadamard relations are recast as a polynomial ideal-membership claim and
discharged on the \emph{same} pipeline.
\emph{(3) The commitment} (\Cref{sec:commitment}): a hash-based $\Ftwo[X]$
commitment whose opening is \emph{bit-sliced}, described next.

\paragraph{Contributions.} Relative to a prime-field arithmetization and to the
integer ($\Z[X]$) pipeline, the construction:
\begin{enumerate}[leftmargin=2em]
  \item \emph{settles the linear and the nonlinear constraints from one
    commitment opening, at one point.} The opening binds, per column, a single
    degree-$<\wbit$ object $a'$ --- the $\eq$-folded $\{0,1\}$ bit slices of the
    committed cells. Both the evaluation $\psia$ (monomial reading, for the
    linear constraints) and the \emph{Lagrange} reading $\psiz$ (for the $\AND$
    relations) are length-$\wbit$ inner products against that \emph{one}
    object, read at the \emph{same} point. So the nonlinear constraints cost no
    second commitment and no second point, and the claim object has fixed size
    $\wbit$, independent of $\log|\K|$. This dual read-off is the feature we
    contribute on top of the underlying binary-field commitment.
  \item \emph{commits one bit per modular addition.} The cell-ring
    representation makes the carry word a free virtual shift, so only the single
    carry-out bit the shift cannot supply is committed --- never a full
    $\wbit$-bit carry word.
  \item \emph{isolates all nonlinearity into column Hadamard products} and
    discharges them, recast as ideal membership, on the same ideal check and
    sumcheck as the linear part.
  \item \emph{runs the whole proof over $\K = \mathrm{GF}(2^{128})$ with no
    prime projection}, the cell ring never widening through encoding.
\end{enumerate}
What is \emph{not} new is worth stating plainly. The per-step full-adder
identity for modular addition and the single-$\AND$ form of $\Ch$ are taken
from Binius64 (\Cref{sec:cmp}), and reading bits as Lagrange coefficients to
turn a coefficient-wise product into a polynomial identity is a standard move.
The contribution is the \emph{composition} --- one bit-sliced opening serving
both the monomial and the Lagrange reading --- and its consequence: the only
nonlinear primitive rides the linear pipeline at no extra commitment cost.

\paragraph{Performance.} \Cref{tab:intro-sweep} previews the headline across a
sweep of chained compressions, both systems at a matched commitment rate
($1/4$, one Apple~M4, CPU). The $\Ftwo[X]$ prover leads at every size and the
lead \emph{widens} with scale --- from near-parity at $10^3$ compressions to
$\sim$$2\times$ at $10^4$ --- while the verifier grows as $\sqrt N$ here against
$O(N)$ for Binius64, an order of magnitude by $8000$ compressions. A second,
structural difference underlies both (\Cref{tab:intro-preproc}): this work is a
\emph{uniform} AIR with $O(1)$ preprocessing, whereas Binius64's non-uniform
word circuit carries an $O(N)$ circuit description that reaches hundreds of
megabytes and whose build \emph{walls} near $16{,}000$ compressions --- where
this work still proves in $\approx 1.2$\,s and stays tractable through
$123{,}366$ compressions ($\nu = 23$). In the other direction, Binius64's
recursive proof is several times smaller; \Cref{sec:cmp} and \Cref{app:meas}
give the full comparison.

\begin{table}[t]
\centering\small
\renewcommand{\arraystretch}{1.2}
\begin{tabular}{r r rr rr}
\toprule
 & & \multicolumn{2}{c}{prover} & \multicolumn{2}{c}{verifier} \\
\cmidrule(lr){3-4}\cmidrule(lr){5-6}
$\approx$comp. & message & this work & Binius64 & this work & Binius64 \\
\midrule
$60$         & $3.8$ KB  & $10.1$ ms          & $33.7$ ms     & $4.8$ ms            & $3.1$ ms \\
$250$        & $16$ KB   & $23.3$ ms          & $41.6$ ms     & $9.3$ ms            & $4.6$ ms \\
$1000$       & $64$ KB   & $75.0$ ms          & $84.1$ ms     & $17.9$ ms           & $14.0$ ms \\
$2000$       & $128$ KB  & $148$ ms           & $155$ ms      & $19.8$ ms           & $27.5$ ms \\
$8000$       & $512$ KB  & $581$ ms           & $1220$ ms     & $45$ ms             & $493$ ms \\
$15{,}420$   & $0.99$ MB & $1190$ ms          & ---$^\dagger$ & $83$ ms             & ---$^\dagger$ \\
\midrule
$30{,}841$ \,($\nu{=}21$)  & $1.97$ MB & $1.35$ s           & ---$^\dagger$ & $117$ ms            & ---$^\dagger$ \\
$61{,}683$ \,($\nu{=}22$)  & $3.95$ MB & $2.72$ s           & ---$^\dagger$ & $224$ ms            & ---$^\dagger$ \\
$123{,}366$ \,($\nu{=}23$) & $7.9$ MB  & $7.0$ s$^\ddagger$ & ---$^\dagger$ & $383$ ms$^\ddagger$ & ---$^\dagger$ \\
\bottomrule
\end{tabular}
\caption{SHA-256 prover and verifier across a sweep of chained compressions at
matched rate $1/4$ on one Apple~M4 (CPU); each compression hashes one
$512$-bit block, so message size $\approx 64\,\text{B}\times$\,compressions.
\emph{Above the rule}: head-to-head
against Binius64 (single runs, \Cref{tab:scaling}); the $\Ftwo[X]$ prover leads
at every size, the lead widening from $\sim$$1.1\times$ at $10^3$ to
$\sim$$2\times$ at $10^4$, and the verifier grows as $\sqrt N$ against $O(N)$,
an order of magnitude by $8000$ compressions. \emph{Below the rule}: beyond
Binius64's build wall, this work alone (deployed-sweep medians,
\Cref{tab:zinc-scaling}), still proving $123{,}366$ compressions ($\nu = 23$).
In return, Binius64's recursive proof is $4$--$7\times$ smaller; ratios are
stabler than the absolute times. $^\dagger$Binius64's $O(N)$ circuit build
exceeds a minute past $\sim$$16{,}000$ compressions. $^\ddagger$This work
memory-bound by swap on the $16$\,GB host at $\nu = 23$; the compute stays
$O(N)$.}
\label{tab:intro-sweep}
\end{table}

\begin{table}[t]
\centering\small
\renewcommand{\arraystretch}{1.2}
\begin{tabular}{r l r}
\toprule
$\approx$comp. & this work & Binius64 \\
\midrule
$16$           & \multirow{6}{*}{$\sim$\,KB \ ($O(1)$)} & $3.7$ MB \\
$64$           &                                        & $14.4$ MB \\
$256$          &                                        & $54.7$ MB \\
$1024$         &                                        & $218$ MB \\
$4096$         &                                        & $873$ MB \\
$\ge 16{,}384$ &                                        & build $> 1$ min \\
\bottomrule
\end{tabular}
\caption{Preprocessing --- the circuit description the system holds, which the
verifier in effect must read --- at matched rate $1/4$. This work is a
\emph{uniform} AIR: one constraint description for any trace height, so
preprocessing is $O(1)$ (a fixed UAIR specification plus the RAA permutation
seeds, kilobyte-scale). Binius64 is a non-uniform word circuit, so its
constraint system plus shift/wiring key-collection grow $O(N)$, the
$\ge 16{,}384$-compression build exceeding a minute. This $O(N)$-vs-$O(1)$ gap
is the root of both Binius64's $O(N)$ verifier and its build wall
(\Cref{tab:intro-sweep}).}
\label{tab:intro-preproc}
\end{table}

\paragraph{Scope and framing.} We present the construction over generic
parameters --- a word width $\wbit$, a binary extension field $\K$, a linear
code --- and record the deployed instantiation ($\wbit = 32$,
$\K = \mathrm{GF}(2^{128})$, seven chained SHA-256 compressions, a rate-$1/4$
$\Ftwo$-RAA code) in remarks and in \Cref{sec:impl}. The commitment is
summarised in \Cref{sec:commitment} and developed in full in a companion note
on the bit-sliced $\Ftwo[X]$ opening; the present paper focuses on how
arithmetization, PIOP, and commitment compose into a single argument. We do not
treat zero-knowledge, and we are explicit that the proof is Brakedown-sized and
that the modular-addition carries are presently bound by a trusted reading
rather than to the commitment (\Cref{sec:soundness}).

\paragraph{Organisation.} \Cref{sec:overview} is a self-contained technical
overview. \Cref{sec:prelim} fixes notation: the cell ring, the field, the trace
model, multilinear extensions, and the two readings. \Cref{sec:arith} develops
the arithmetization. \Cref{sec:piop} gives the binary-field PIOP for the linear
constraints and \Cref{sec:hadamard} the Hadamard discharge. \Cref{sec:commitment}
summarises the bit-sliced commitment. \Cref{sec:soundness} composes the
soundness of the layers under Fiat--Shamir, and \Cref{sec:impl} reports the
deployed instantiation and its costs. \Cref{sec:cmp} positions the construction
against Binius64.
